\section{Chellas's standard systems: source constructions}

\begin{annexitem}{annex:p8.chellas.normal-interface}{thm:p08-chellas-normal-interface}{Normal closure, theorem-only proofs and the syntax bridge}
\begin{proof}
The empty-list instance of normal closure is theorem-only necessitation.
The one-antecedent instance boxes a theorem implication. Applied to
\(((A\to B)\land A)\to B\), normal closure yields
\((\Box(A\to B)\land\Box A)\to\Box B\); propositional reasoning
gives K. Conversely, K and necessitation imply that a theorem implication
can be boxed and distributed. Necessitate a theorem
\(A_1\to(A_2\to\cdots(A_r\to B)\cdots)\); repeated K and
modus ponens with \(\Box A_1,\ldots,\Box A_r\) yield \(\Box B\).
Discharge these premises propositionally to get the source's boxed
consequence. For \(r=0\), necessitate the original theorem itself.

The deduction theorem for finite premise derivations is induction on
proof lines. A premise equal to the discharged \(A\) gives the
tautology \(A\to A\). Every other premise or axiom \(B\) gives
\(A\to B\) using \(B\to(A\to B)\). The modus-ponens case
uses \((A\to(B\to C))\to((A\to B)\to(A\to C))\).
A necessitated line is already a theorem from a separate no-premise
proof, so it is handled like an axiom, not boxed again under the premise.
Repeated deduction gives exactly the finite-conjunction implication
defining source deducibility. Conversely, that implication and its
finitely many premises give its consequent by propositional inference.
Thus the two consequence definitions agree for arbitrary premise sets.

A theorem \(A\leftrightarrow B\) implies both directions. Boxing
them and using K gives \(\Box A\leftrightarrow\Box B\).
Negation and \(\mathrm{Df}\Diamond\) give the corresponding diamond
equivalence. Structural induction now permits replacement at any
chosen occurrence: at Boolean nodes use the corresponding tautological
congruence, and at modal nodes use these two congruences. If the whole
sentence is the replaced occurrence, use the original equivalence;
unchanged branches use the identity equivalence. This also covers no
replacements and several separately chosen occurrences.

Define a translation \(e\) recursively, mapping \(P_n\) to \(p_n\),
preserving the primitive Boolean connectives and box, expanding
\(\topv\) and biconditional in their usual Boolean definitions,
and sending \(\Diamond A\) to \(\neg\Box\neg e(A)\).
In source syntax, interpret its output back by replacing \(p_n\)
with \(P_n\). Structural induction with replacement and
\(\mathrm{Df}\Diamond\) proves \(A\leftrightarrow e(A)\)
under that interpretation. The same recursion gives truth equivalence
in every standard model, including the diamond case by the quantified
duality of its two truth clauses.

To prove theoremhood equivalence, translate every source proof line.
Propositional tautologies stay tautological, each additional schema
translates to the corresponding sheet schema, and K, modus ponens and
theorem-only necessitation translate to the same operations. The
\(\mathrm{Df}\Diamond\) instances translate to Boolean identities.
In the reverse direction, interpret each sheet proof as a source proof
and replace its translated last sentence by its theorem-equivalent
source sentence. Apply the same argument to the finite-conjunction
implication to obtain premise-consequence equivalence. Normal closure
of each selected schema set is therefore exactly its sheet theory.
\end{proof}
\end{annexitem}

\begin{annexitem}{annex:p8.chellas.correspondence}{thm:p08-chellas-correspondence}{Exact relational correspondences and local soundness}
\begin{proof}
Seriality makes a true box have a witnessing successor, proving D;
reflexivity makes its current world a successor, proving T. Symmetry
makes the current true-A world available from every successor, proving
B. Transitivity passes truth of A at every current successor to every
successor's successors, proving 4. For 5, if \(\alpha R\beta\)
and A is true at \(\beta\), Euclideanness makes every other successor
\(\gamma\) of \(\alpha\) access \(\beta\), so each such
\(\gamma\) satisfies \(\Diamond A\). For G, suppose
\(\alpha R\beta\) and \(\Box A\) holds at \(\beta\).
For any \(\alpha R\gamma\), convergence provides
\(\beta R\delta\) and \(\gamma R\delta\).
A is true at \(\delta\), so \(\Diamond A\) holds at \(\gamma\).

For the converses use arbitrary valuations of one atom. At a world
without successors, \(\Box\topv\to\Diamond\topv\) fails.
If \(\alpha R\alpha\) fails, make P true at every world other
than \(\alpha\); \(\Box P\to P\) fails there.
If \(\alpha R\beta\) but not \(\beta R\alpha\), put
\(P=\{\alpha\}\); P is true at \(\alpha\), while its successor
\(\beta\) falsifies \(\Diamond P\), refuting B.
If \(\alpha R\beta R\gamma\) but not \(\alpha R\gamma\),
make P true at exactly the successors of \(\alpha\). This makes
\(\Box P\) true at \(\alpha\) and false at its successor
\(\beta\), refuting 4. If \(\alpha R\beta,\alpha R\gamma\)
but not \(\beta R\gamma\), put \(P=\{\gamma\}\).
Then \(\Diamond P\) holds at \(\alpha\) but fails at the
successor \(\beta\), refuting 5. Finally, if two successors
\(\beta,\gamma\) of \(\alpha\) have no common successor,
make P true precisely at the successors of \(\beta\).
Then \(\Box P\) holds at \(\beta\), while \(\Diamond P\)
fails at \(\gamma\), refuting G at \(\alpha\).
These arguments include coincident worlds and empty successor sets.

K is valid on every frame: every successor satisfying A and A-implies-B
satisfies B. Tautologies and modus ponens preserve truth at a world.
A no-premise theorem true at every world can be necessitated, since it
is true at every successor. Induction on no-premise proof trees proves
theoremhood soundness, and then induction on local premise derivations
proves consequence soundness. Equivalently, the source finite theorem
implication is valid everywhere and its true antecedents imply its
consequent at the selected world. Each additional schema has just been
checked on its exact frame class.

If R is reflexive and Euclidean, \(\alpha R\beta\) together with
\(\alpha R\alpha\) gives \(\beta R\alpha\), proving symmetry.
If \(\alpha R\beta R\gamma\), symmetry gives
\(\beta R\alpha\); Euclideanness at \(\beta\) yields
\(\alpha R\gamma\). Conversely, reflexivity, symmetry and
transitivity imply Euclideanness by reversing the first edge and
composing. This proves the equivalence-frame convention.
\end{proof}
\end{annexitem}

\begin{annexitem}{annex:p8.chellas.selected-inclusions}{thm:p08-chellas-selected-inclusions}{Schema derivations and every noninclusion}
\begin{proof}
Adding schemas gives the obvious inclusions. To derive D from T, duality
and replacement give \(A\to\Diamond A\); combine this with
\(\Box A\to A\). From T and 5, combine
\(A\to\Diamond A\to\Box\Diamond A\) to get B.
In a B-system, a theorem \(\Diamond A\to B\) implies
\(A\to\Box B\): box the theorem, distribute, and combine it
with \(A\to\Box\Diamond A\). Apply this operation to the
dual 4 theorem \(\Diamond\Diamond A\to\Diamond A\) to get
5, and to the dual 5 theorem \(\Diamond\Box A\to\Box A\)
to get 4. Thus B plus either 4 or 5 gives the other.
In a D/B/4 system, combine
\(\Box A\to\Box\Box A\to\Diamond\Box A\to A\),
using 4, the D instance for \(\Box A\), and the dual B theorem.
This gives T and establishes the advertised S5 axiomatizations.
In a B-system, box the dual B theorem
\(\Diamond\Box A\to A\), and compose with the B instance for
\(\Diamond\Box A\); the result is
\(\Diamond\Box A\to\Box\Diamond A\), namely G.
Therefore S5 includes S4.2.

Here are complete separating frames; unlisted edges are absent and
unmentioned atoms can be false everywhere.
\begin{enumerate}
\item A single world with no edge validates K, KB and K4 and falsifies
 D (take \(\topv\)), hence also T.
\item On \(\{a,b\}\), the relation \(\{(a,b),(b,b)\}\) is
 serial and transitive. With \(P_0=\{b\}\), T fails at a.
 With \(P_0=\{a\}\), B fails at a. Thus neither KD nor K4
 supplies either reflexivity or symmetry.
\item On \(\{a,b\}\), the relation \(\{(a,b),(b,a)\}\) is
 symmetric and serial. With \(P_0=\{b\}\), T fails at a.
 With \(P_0=\{b\}\), \(\Box P_0\) is true at a but false
 at b, refuting 4 at a. This separates KB from KT and K4.
\item On \(\{a,b,c\}\), take all reflexive edges and both
 directions between a and b and between b and c. With
 \(P_0=\{a,b\}\), 4 fails at a. This KTB frame therefore
 separates KTB from every selected system containing 4.
\item On \(\{a,b\}\), take reflexive edges and \((a,b)\).
 With \(P_0=\{a\}\), B fails at a. This frame is S4.2:
 b is a common successor of every pair of worlds. It separates
 S4 and S4.2 from KB, KTB and S5.
\item On \(\{a,b,c\}\), take reflexive edges and
 \((a,b),(a,c)\). This S4 fork is not directed. With
 \(P_0=\{b\}\), \(\Diamond\Box P_0\) holds at a, but
 \(\Box\Diamond P_0\) fails there because c has no P-successor.
 It separates S4, hence also KT and K4, from S4.2.
\end{enumerate}
The first six cases also prove every declared edge strict. For completeness
of the noninclusion list, a selected system is determined syntactically by
the following required characteristic features: K has none; KD has D;
KT has T; KB has B; KTB has T/B; K4 has 4; S4 has T/4;
S4.2 has T/4/G; S5 has T/B/4/5. The schema derivations above account
for all redundant features (T entails D, B entails G, and S5 contains
the rest). A pair outside the generated inclusion order requires in
one direction a missing T, B, 4, D or G. The corresponding numbered
model satisfies the smaller system and falsifies that missing feature:
isolated worlds for D; the serial arrow or symmetric two-cycle for T;
the convergent chain for B; the reflexive symmetric path for 4;
and the reflexive fork for G. Consequently no additional inclusion
among the nine systems is possible. This argument is only about this
finite selection.
\end{proof}
\end{annexitem}

\begin{annexitem}{annex:p8.chellas.source-canonical}{thm:p08-chellas-source-canonical}{Construction rather than an assumed canonical model}
\begin{proof}
Finitary deducibility has finite support by definition. If a consistent
X becomes inconsistent upon adding A and also upon adding \(\neg A\),
the deduction theorem and propositional reasoning make X inconsistent.
Enumerate all sentences. Starting with X, add the next sentence when
consistent, otherwise add its negation. Every stage is consistent.
An inconsistent union would have a finite inconsistent subset, all of
whose members occur in one stage (take the maximum of their least
stage indices), a contradiction. The union decides every sentence
and has no consistent proper extension, hence is maximal. This
classical recursion is an existence construction, not a decision
algorithm for consistency of arbitrary X.

In any maximal set \(\alpha\), derivability is membership. Indeed,
if \(\alpha\) derives A but omits A, maximality makes
\(\alpha\cup\{A\}\) inconsistent; substitute the derivation of A
for each added premise to contradict consistency of \(\alpha\).
Thus every theorem belongs to it and it is closed under propositional
inference. It contains exactly one of A and \(\neg A\), excludes
\(\bottom\), includes \(\topv\), contains a conjunction exactly
when it contains both conjuncts, contains a disjunction exactly when
it contains at least one disjunct, and contains \(A\to B\) exactly
when A is absent or B present. For conjunction, use the projections
and the introduction tautology. For disjunction, if neither disjunct
is present, both negations are, contradicting the disjunction. For
implication, modus ponens proves the necessary direction, while
\(\neg A\to(A\to B)\) and \(B\to(A\to B)\) prove sufficiency.
Biconditional membership follows from its equivalence to the two
implications. These establish every Boolean membership clause.

Let \(U(\alpha)=\{B:\Box B\in\alpha\}\). If
\(\neg\Box A\in\alpha\), then
\(U(\alpha)\cup\{\neg A\}\) is consistent. Otherwise a finite
list \(B_i\in U(\alpha)\) would deduce A, and normal boxed
consequence would put \(\Box A\) in \(\alpha\), a contradiction.
This includes the empty list by theorem-only necessitation.
Maximally extend the consistent set to \(\beta\). Then
\(\alpha R_\Sigma\beta\) and \(\neg A\in\beta\).
Duality also gives a successor containing A exactly when
\(\Diamond A\in\alpha\). Further,
\(U(\alpha)\subseteq\beta\) is equivalent to
\(\{\Diamond A:A\in\beta\}\subseteq\alpha\): failure of
either inclusion supplies, through maximal negation and duality,
a counterexample to the other. Explicitly, \(A\in\beta\) and
\(\Diamond A\notin\alpha\) imply \(\Box\neg A\in\alpha\),
contradicting the first inclusion; conversely, \(\Box A\in\alpha\)
and \(A\notin\beta\) imply \(\neg A\in\beta\), so the second
inclusion supplies \(\Diamond\neg A\in\alpha\), a contradiction.

The empty theory is consistent and hence has a maximal extension,
so the constructed world set is nonempty. Truth by structural induction
now agrees with membership. For atoms this is the valuation; constants
and every Boolean case are the preceding membership clauses. A member
\(\Box A\) makes all successors contain A, while a missing box has
the successor containing \(\neg A\) just constructed. Induction
therefore proves both directions of the box case. A diamond member has
a successor containing its matrix, while a missing diamond has no such
successor by the dual accessibility equivalence. This proves both
directions of the primitive diamond case as well.

A theorem belongs to every maximal set and is thus true everywhere
under the canonical valuation. A non-theorem A makes \(\{\neg A\}\)
consistent and has a maximal extension falsifying A. Therefore the
single constructed model determines theoremhood. This does not yet
assert validity under all valuations on its frame: that requires the
separate frame-property and soundness arguments.
\end{proof}
\end{annexitem}

\begin{annexitem}{annex:p8.chellas.kb-completeness}{thm:p08-chellas-kb-completeness}{The symmetry-only canonical frame}
\begin{proof}
The one-world reflexive model validates K and B and falsifies
\(\bottom\), so KB is consistent. Construct its proper canonical
model by the preceding lemma. If \(\alpha R\beta\) and
\(A\in\alpha\), B gives \(\Box\Diamond A\in\alpha\),
so \(\Diamond A\in\beta\). The dual accessibility equivalence
therefore gives \(\beta R\alpha\). Only B, not T, was used.

If \(\Gamma\not\vdash_{KB}A\),
\(\Gamma\cup\{\neg A\}\) is consistent by finite support and
deduction. Its maximal extension is a canonical world satisfying every
member of \(\Gamma\) and falsifying A. The frame is symmetric by
the argument above. This proves strong local completeness; soundness
is Annex~\ref{annex:p8.chellas.correspondence}.
The two-world symmetric two-cycle refutes T and 4, and the isolated
world refutes D, as given in
Annex~\ref{annex:p8.chellas.selected-inclusions}. Hence none of
reflexivity, transitivity or seriality can be added to the conclusion
as a required property of all KB frames. In particular KB is not KTB.
\end{proof}
\end{annexitem}

\begin{annexitem}{annex:p8.chellas.g-determination}{thm:p08-chellas-g-determination}{All indices, path factorization and canonical convergence}
\begin{proof}
Induction on j gives the path truth clauses: \(\Box^j A\) is
true at \(\alpha\) exactly when A is true at all j-step successors;
\(\Diamond^j A\) is true exactly when some j-step successor
satisfies A. For j=0 both claims are the identity clause. For j+1,
apply the one-step modal clause and then the induction hypothesis
at each one-step successor; the resulting two-stage quantification
is precisely composition of path relations. Consequently the
four-index frame condition validates the schema: a k-step witness
\(\beta\) for \(\Box^l A\), together with any m-step successor
\(\gamma\), has a common \(\delta\) l steps from \(\beta\)
and n steps from \(\gamma\). A holds there, witnessing
\(\Diamond^n A\) at \(\gamma\). Zero indices require no exception.
The one-world reflexive frame satisfies the condition for every
choice of indices and falsifies \(\bottom\), so the least normal
schema extension is consistent.

We prove the needed path lemma for the proper canonical relation of
any consistent normal \(\Sigma\):
\[
 \alpha R_\Sigma^j\beta
 \quad\Longleftrightarrow\quad
 U_j(\alpha)\subseteq\beta,
 \qquad U_j(\alpha)=\{A:\Box^j A\in\alpha\}.
\]
For j=0, \(U_0(\alpha)=\alpha\), and containment between maximal
consistent sets is equality: a proper extension cannot be consistent.
For j=1 this is the definition of the proper relation. Repeated
normal closure also gives the iterated boxed-consequence operation:
from a theorem \(\bigwedge A_i\to B\), derive
\(\bigwedge\Box^j A_i\to\Box^j B\) by induction on j, with
the empty list at every stage included.

The corresponding dual containment statement is
\[
 U_j(\alpha)\subseteq\beta
 \quad\Longleftrightarrow\quad
 \{\Diamond^j A:A\in\beta\}\subseteq\alpha.
\]
Its proof repeats the maximal-negation argument for j=1, using
\(\Diamond^j A\leftrightarrow\neg\Box^j\neg A\), which follows
by induction on j, modal congruence and duality. This identity is
also valid when j=0.

For the inductive path factorization, suppose
\(U_{j+1}(\alpha)\subseteq\beta\). The set
\[
 U_1(\alpha)\cup\{\Diamond^j C:C\in\beta\}
\]
is consistent. Otherwise finite lists \(B_i\in U_1(\alpha)\)
and \(C_t\in\beta\) give a theorem
\(\bigwedge B_i\to\neg\bigwedge\Diamond^j C_t\).
Put \(C=\bigwedge C_t\), using \(\topv\) for the empty list.
Monotonicity gives
\(\Diamond^j C\to\bigwedge\Diamond^j C_t\), including the
empty-conjunction case. Thus
\(\bigwedge B_i\to\Box^j\neg C\) is a theorem. Box this finite
consequence to put \(\Box^{j+1}\neg C\) in \(\alpha\).
The assumed inclusion puts \(\neg C\) in \(\beta\), whereas
\(C\in\beta\), a contradiction. Maximally extend the union to
\(\gamma\). Then \(\alpha R\gamma\), and the dual containment
statement gives \(U_j(\gamma)\subseteq\beta\). The induction
hypothesis gives \(\gamma R^j\beta\), hence
\(\alpha R^{j+1}\beta\). Conversely, if
\(\alpha R\gamma R^j\beta\) and \(\Box^{j+1}A\in\alpha\),
then \(\Box^j A\in\gamma\), so \(A\in\beta\).
This proves both directions of factorization and the path lemma.

Now assume \(\Sigma\) contains \(G^{k,l,m,n}\), and suppose
\(\alpha R^k\beta\) and \(\alpha R^m\gamma\). The set
\(U_l(\beta)\cup U_n(\gamma)\) is consistent. If not,
finite \(B_i\in U_l(\beta)\), \(C_t\in U_n(\gamma)\)
would yield \(\bigwedge B_i\to\neg\bigwedge C_t\).
Writing \(C=\bigwedge C_t\), the iterated boxed-consequence
operation with index l gives \(\Box^l\neg C\in\beta\).
The k-step dual containment statement puts
\(\Diamond^k\Box^l\neg C\) in \(\alpha\). The G instance puts
\(\Box^m\Diamond^n\neg C\) there, and the m-step path lemma
puts \(\Diamond^n\neg C\) in \(\gamma\). But the C-list
also puts \(\Box^n C\) in \(\gamma\), contradicting iterated
duality. Empty lists use the boxed theorem \(\Box^l\topv\) or
\(\Box^n\topv\); the argument also covers l=0 or n=0.
A maximal extension \(\delta\) of this consistent union gives
\(\beta R^l\delta\) and \(\gamma R^n\delta\), exactly
the required canonical frame condition.

For the least normal extension, soundness follows from the first
paragraph. If \(\Gamma\not\vdash_\Sigma A\), maximally extend
\(\Gamma\cup\{\neg A\}\) and use the canonical truth lemma.
The preceding frame argument supplies a countermodel in the stated
class for the entire original premise set. This proves strong local
completeness, and theoremhood determination is its empty-premise case.
For a stronger normal extension, the same canonical frame calculation
still applies, but soundness on every frame with only the one condition
is not implied. This is why the theorem keeps those two claims distinct.
\end{proof}
\end{annexitem}

\begin{annexitem}{annex:p8.chellas.s42-completeness}{thm:p08-chellas-s42-completeness}{Canonical local convergence and the generated source class}
\begin{proof}
The S4.2 proper canonical frame is reflexive by T: a boxed matrix
belongs to its own world. It is transitive by 4: if
\(\alpha R\beta R\gamma\) and \(\Box A\in\alpha\), then
\(\Box\Box A\in\alpha\), so \(\Box A\in\beta\), then
\(A\in\gamma\). The previous annex item with k=l=m=n=1
makes it locally directed by G. These calculations apply simultaneously
in the same canonical model, not in three unrelated models.
The one-world reflexive model validates all its axioms, hence the
system is consistent. Soundness and the full canonical truth construction
therefore prove strong local completeness for arbitrary premises.

Restrict a canonical countermodel to the worlds generated by its
evaluation world \(\alpha\). Each retained world is a successor
of \(\alpha\), since the frame is reflexive and transitive.
Any two retained worlds \(\beta,\gamma\) have a common successor
\(\delta\) by directedness at \(\alpha\), and \(\delta\)
is retained by transitivity. Thus the restriction is globally
convergent, reflexive and transitive. No successor of a retained world
is discarded. Structural induction preserves all formula truth:
atoms and Booleans are unchanged, and a modal sentence sees exactly
the same successors with the same matrix truth. The distinguished
world therefore retains every premise and still falsifies the conclusion.
This proves the stronger class restriction. Every globally convergent
frame is locally directed, so soundness is unchanged.
\end{proof}
\end{annexitem}

\begin{annexitem}{annex:p8.chellas.generated-truth}{thm:p08-chellas-generated-truth}{All formula cases in generated restrictions}
\begin{proof}
The generated set contains its root by the zero path. If a retained
\(\beta\) has an R-successor \(\gamma\), append that edge to
its finite path from the root to see that \(\gamma\) is retained.
The original and restricted successor sets are consequently identical
at every retained world. Structural induction at all retained worlds
proves truth preservation. An atom has the same membership in its
restricted valuation; constants have the same fixed truth value;
negation reverses equal truth values; conjunction, disjunction,
conditional and biconditional apply the same Boolean operation to
equal operand values. For box, universal quantification runs over
the same successors and the induction hypothesis agrees there.
For diamond, the same existential witnesses agree in both models.

If A is true everywhere in an original model, it is true in every
generated restriction. Conversely, if it is false at a world,
the model generated by that very world still falsifies it there.
This proves the class-validity equality without requiring a model to
be generated by every one of its worlds.
In an equivalence frame, the worlds generated by \(\alpha\) are
exactly its equivalence class, on which R is universal. Universal
frames are themselves equivalence frames and generated from any
world. Strong S5 completeness together with this truth preservation
therefore gives determination by universal models. Filtering such
a generated model through its finitely many subsentences preserves
universality by original-edge containment, gives at most \(2^s\)
worlds and preserves its falsified sentence by the filtration truth
lemma below. This proves finite universal determination as well.
The argument uses no filtration of an arbitrary equivalence relation
that is asserted, without proof, to remain an equivalence relation.
\end{proof}
\end{annexitem}

\begin{annexitem}{annex:p8.chellas.filtration-truth}{thm:p08-chellas-filtration-truth}{Every primitive case and both quotient relations}
\begin{proof}
Agreement on all \(\Gamma\)-sentences is reflexive, symmetric and
transitive, so its classes partition W. Since every matrix of a modal
sentence in \(\Gamma\) is also in \(\Gamma\), every defining
constraint on \(R^*\) is invariant under changing representatives.
The finest relation contains all original edges by definition. If some
representatives \(\alpha' R\beta'\) witness an edge, a true
\(\Box A\in\Gamma\) at \(\alpha\) is true at \(\alpha'\);
A is true at \(\beta'\), hence at \(\beta\). For the diamond
constraint, if A is true at \(\beta\), it is true at \(\beta'\),
so \(\Diamond A\) is true at \(\alpha'\), hence at \(\alpha\).
Thus the finest relation is a filtration. The coarsest relation
satisfies both modal requirements by its definition; every original
edge meets them by the original truth clauses, so it too is a filtration.
The same reasoning applies to any intermediate relation.

Induct on A in \(\Gamma\), simultaneously at every world.
For atoms, \([\alpha]\in[P_n]\) is equivalent to
\(\alpha\in P_n\), because atom membership is constant on a
class. Constants are unchanged. Negation and each of conjunction,
disjunction, implication and biconditional use the same Boolean
truth operation on operands in \(\Gamma\), so the induction
hypothesis proves the required equivalence in each case.
If \(\Box B\) is true in the original model, the box constraint
makes B true at every quotient successor; induction gives a true
quotient box. If the original box is false, choose an original
successor falsifying B; its quotient edge persists, and induction
still falsifies B there. If an original diamond is true, its original
witness edge persists and its matrix stays true, making the quotient
diamond true. Conversely, a quotient successor satisfying B corresponds,
by induction, to an original world satisfying B; the diamond constraint
makes \(\Diamond B\) true at the original current world. This
proves the primitive diamond case without assuming its dual expansion
lies in \(\Gamma\).

The Boolean closure adds only Boolean combinations of sentences
whose truth has just been preserved. A second structural induction
on those combinations proves its extended truth equivalence.
This does not preserve arbitrary additional modalizations.
The map sending a class to its truth pattern on s representatives
is injective: equal patterns give agreement on every
\(\Gamma\)-sentence, hence the same class. There are \(2^s\)
patterns. This proves both the finite syntactic and the logically
finite bounds, even when \(\Gamma\) itself is infinite.
Original seriality supplies a successor edge for every class;
original reflexivity supplies every self-edge; and original
universality supplies every pair of classes. These prove the three
automatic preservation claims.
\end{proof}
\end{annexitem}

\begin{annexitem}{annex:p8.chellas.paired-filtrations}{thm:p08-chellas-paired-filtrations}{Symmetric, transitive and Euclidean refinements}
\begin{proof}
Use the abbreviations \(B_A(\alpha)\) and \(D_A(\alpha)\)
for original truth of \(\Box A\) and \(\Diamond A\), respectively,
only when the corresponding modal sentence is in \(\Gamma\).
The c-conditions are class-invariant by their definition.
In every construction c1 guarantees the two modal filtration constraints.
It therefore remains to check the relation properties and that every
original edge satisfies the additional conditions.

For c1/c2 symmetry is immediate from the inclusion of both the
forward and reverse c1 constraints. If the original model is symmetric,
an original edge and its reverse both satisfy c1, hence its quotient
edge belongs to the constructed relation.

For c1/c3 transitivity, suppose the constraints hold for
\((\alpha,\beta)\) and \((\beta,\gamma)\).
A box true at \(\alpha\) propagates by c3 to \(\beta\),
then by c1 to its matrix at \(\gamma\); a matrix true at
\(\gamma\) propagates by c1 to a diamond at \(\beta\),
then by c3 to a diamond at \(\alpha\). These prove c1 for
\((\alpha,\gamma)\). A box true at \(\alpha\) propagates
twice by c3 to a box at \(\gamma\), and a diamond true at
\(\gamma\) propagates twice backwards to a diamond at
\(\alpha\), proving c3. On an original transitive edge, a box
at the source remains true at the target because every target
successor is a source successor; a diamond at the target gives a
source witness by transitivity. Thus original edges satisfy c3.

For c1/c2/c3/c4, both c1 and c3 are accompanied by their reverse
constraints, so the constructed relation is symmetric. It is the
intersection of the transitive c1/c3 relation and its converse;
the converse of a transitive relation is transitive, and intersection
preserves transitivity. An original symmetric transitive edge satisfies
c1/c3 in both directions by the preceding arguments, giving every
required original quotient edge. This proves the symmetric transitive
construction with all its constraints explicitly present.

For c1/c3/c4, transitivity follows by the c1/c3 argument and
transitivity of equality of the modal truth records imposed by
c3/c4. To check Euclideanness, suppose
\([\alpha]R^*[\beta]\) and \([\alpha]R^*[\gamma]\).
The box and diamond truth records at all three worlds agree, by c3/c4.
If \(B_A(\beta)\), then \(B_A(\alpha)\), and c1 on the
second edge gives A at \(\gamma\). If A holds at \(\gamma\),
c1 gives \(D_A(\alpha)\), whose agreed record gives
\(D_A(\beta)\). These prove c1 for \((\beta,\gamma)\);
the agreed records give c3/c4 there too, so
\([\beta]R^*[\gamma]\). On an original edge of a transitive
Euclidean model, transitivity gives c3. Euclideanness gives c4:
a box true at the target constrains every source successor because
the target accesses all of them; and any source diamond witness
is a successor of the target. Thus original edges satisfy all three
conditions. Every constructed relation is a filtration, and automatic
seriality/reflexivity preservation follows from the preceding annex item.
\end{proof}
\end{annexitem}

\begin{annexitem}{annex:p8.chellas.modal-closure-filtration}{thm:p08-chellas-modal-closure-filtration}{Coarsest modal-closure filtrations and logical finiteness}
\begin{proof}
All the extra modal sentences used below belong to \(\Gamma\)
by modal closure, including prefixed negations. Suppose first that
the original model is symmetric and \([\alpha]R^*[\beta]\).
If \(\Box A\) is true at \(\beta\) and A false at \(\alpha\),
B in the original frame makes \(\Box\Diamond\neg A\) true
at \(\alpha\). Its c1 constraint puts \(\Diamond\neg A\)
at \(\beta\), contradicting the true box there. Hence A holds
at \(\alpha\). If A holds at \(\alpha\), B puts
\(\Box\Diamond A\) there and c1 gives \(\Diamond A\)
at \(\beta\). These are the two reverse c1 constraints,
so the coarsest relation is symmetric.

For transitivity, suppose \([\alpha]R^*[\beta]R^*[\gamma]\)
and the original relation is transitive. If \(\Box A\) holds
at \(\alpha\), 4 gives \(\Box\Box A\) there. Coarsest
constraints give \(\Box A\) at \(\beta\), then A at
\(\gamma\). If A holds at \(\gamma\), the second edge
gives \(\Diamond A\) at \(\beta\), and the first edge
gives \(\Diamond\Diamond A\) at \(\alpha\).
Transitivity collapses it to \(\Diamond A\), by composing
its two witness edges. Hence both c1 constraints hold for
\((\alpha,\gamma)\), proving quotient transitivity.

For Euclideanness, let \([\alpha]R^*[\beta]\) and
\([\alpha]R^*[\gamma]\). If \(\Box A\) holds at
\(\beta\), modal closure permits the first edge to yield
\(\Diamond\Box A\) at \(\alpha\). In an Euclidean frame
this implies \(\Box A\): its witness accesses every successor
of \(\alpha\). The second edge then gives A at \(\gamma\).
If A holds at \(\gamma\), the second edge gives
\(\Diamond A\) at \(\alpha\). Euclideanness gives
\(\Box\Diamond A\) there, and the first edge gives
\(\Diamond A\) at \(\beta\). Again both c1 constraints
hold for \((\beta,\gamma)\), proving the result.

For logical finiteness in an Euclidean model, fix a world \(\alpha\).
If it has no successor, any nonempty box/diamond word has the constant
truth value determined by its first operator. Suppose instead that
\(\alpha R\beta\). Every successor \(\beta\) is reflexive,
by Euclideanness with the repeated successor. Any two successors
of \(\alpha\) access each other. For such two worlds
\(\beta,\gamma\), their successor sets coincide: if
\(\beta R\delta\), the edges \(\beta R\gamma\) and
\(\beta R\delta\) give \(\gamma R\delta\), and the
reverse follows similarly. Call this common set X. It is nonempty,
contains every successor of \(\alpha\), and is a universal
closed cluster. Indeed, any \(\delta\in X\) accesses
\(\beta\) and every member of X by Euclideanness at
\(\beta\); conversely, if \(\delta R\epsilon\), its edge
back to \(\beta\) makes \(\beta R\epsilon\), so
\(\epsilon\in X\).

At every world in that cluster, a one-step box or diamond of A
therefore has a fixed truth value independent of the cluster world.
For any three box/diamond operators, the middle operation quantifies
a constant value over the same nonempty X and can be removed.
Thus every length-three word is model-equivalent to its first and
last operators; the no-successor case agrees as well. Iterating
this reduction leaves only identity, box, diamond, double box,
double diamond, box-diamond or diamond-box. Push negations to the
outside using modal duality and eliminate pairs of negations.
There are at most fourteen resulting modalities, counting their
negations. This argument proves their equivalence at every world
simultaneously, not merely at one selected root. Hence modal closure
of s base sentences has at most \(14s\) truth-set representatives.
The logically finite size argument in
Annex~\ref{annex:p8.chellas.filtration-truth} gives a finite quotient.
This auxiliary construction needs no prior finite-model or
decidability assumption.
\end{proof}
\end{annexitem}

\begin{annexitem}{annex:p8.chellas.filtration-boundaries}{thm:p08-chellas-filtration-boundaries}{Complete failure examples for uncontrolled filtrations}
\begin{proof}
For failure of symmetry in an arbitrary filtration, take two worlds
a,b with no original edges, and \(P_0=\{a\}\),
\(\Gamma=\{P_0\}\). Its classes are singletons. Because there
are no modal sentences in \(\Gamma\), either modal constraint
is vacuous and any quotient relation is a filtration. Choose only
\([a]R^*[b]\), with no reverse edge. The original empty relation
is symmetric, while this filtration is not. In contrast, a finest
filtration of a symmetric model is symmetric because a witnessed
original edge always has its witnessed reverse.

For a finest filtration failing transitivity, take four worlds
a,b,c,d and exactly the edges \((a,b),(c,d)\), with
\(P_0=\{b,c\}\) and \(\Gamma=\{P_0\}\). The original
relation is transitive vacuously: neither target has a successor.
The classes are \([a]=[d]\) and \([b]=[c]\), and the finest
relation has exactly the two opposite edges between them, but no
self-edge. Composing either pair refutes transitivity.

For failure of local directedness, use four worlds a,b,c,d,
all reflexive edges, and only the additional edges
\((a,b),(c,d)\). This relation is reflexive, transitive and
locally directed: each two-element chain has its own terminal
world as a common upper world. Set
\(P_0=\{b\},P_1=\{d\}\) and
\(\Gamma=\{P_0,P_1\}\). The finest quotient has the three
classes \(\{a,c\},\{b\},\{d\}\), all reflexive edges,
and arrows from the first class to each of the other two.
The latter classes have no common successor. The original model
was not globally convergent, precisely the condition supplied by
generated S4.2 restriction in the finite-model argument.

For the stronger Euclidean obstruction, use five distinct worlds
\(\alpha,\beta,\gamma,\delta,\epsilon\). The relation has
\((\alpha,\beta),(\beta,\beta),(\gamma,\delta)\), and all
four pairs on \(\{\delta,\epsilon\}\). Put
\(P_0=\{\beta,\delta\}\), and
\(\Gamma=\{P_0,\Box P_0\}\). Every original successor
set is empty or a subset of a universal cluster whose members
access one another, so the relation is Euclidean. The truth
records \((P_0,\Box P_0)\) are
\[
 \alpha:(0,1),\quad\beta:(1,1),\quad
 \gamma:(0,1),\quad\delta:(1,0),\quad\epsilon:(0,0).
\]
Only \(\alpha\) and \(\gamma\) merge. Any filtration must
therefore contain edges
\([\alpha]\to[\beta]\), \([\alpha]\to[\delta]\),
and \([\delta]\to[\epsilon]\). If its relation were
Euclidean, the first two edges would force both
\([\beta]\to[\delta]\) and \([\delta]\to[\beta]\).
The last forced edge and \([\delta]\to[\epsilon]\) would
force \([\beta]\to[\epsilon]\). But the true
\(\Box P_0\) at \(\beta\) and false P at \(\epsilon\)
prohibit that edge by the box filtration constraint. Hence there
is no Euclidean filtration through this particular \(\Gamma\).
Modal closure avoids exactly this forced merging, as proved above.
\end{proof}
\end{annexitem}

\begin{annexitem}{annex:p8.chellas.finite-models}{thm:p08-chellas-finite-models}{Each selected row and its exact finite bound}
\begin{proof}
Let A be a non-theorem of the selected system and put
\(\Gamma=\operatorname{Sn}(A)\), with s members.
Strong local completeness in the matching canonical theorem supplies
a pointed model of the stated class falsifying A. The primitive-language
bridge preserves truth and theoremhood, so this can be read directly
as a source model. Filter this model through \(\Gamma\), using the
following construction in each case.

For K choose the finest relation, with no extra property needed.
For KD and KT use the same relation; automatic preservation retains
seriality and reflexivity, respectively. For KB use c1/c2 to retain
symmetry. For KTB use c1/c2 as well; reflexivity persists automatically.
For K4 use c1/c3 to retain transitivity. For S4 use c1/c3 and
automatic reflexivity. For S5 one may use c1/c2/c3/c4 to retain
symmetry and transitivity, with reflexivity automatic; an alternative
is its generated universal restriction, any filtration of which
remains universal. Every one of these constructions preserves A
and has at most \(2^s\) worlds.

For S4.2 first restrict to the model generated by the falsifying
world. Annex~\ref{annex:p8.chellas.s42-completeness} proves that it
is globally convergent and that it still falsifies A. Its c1/c3
filtration is reflexive and transitive. For any quotient classes
\([\beta],[\gamma]\), take representatives \(\beta,\gamma\)
in that generated model. Global convergence supplies a retained
\(\delta\) with \(\beta R\delta\) and
\(\gamma R\delta\). Every filtration contains the corresponding
two edges, so \([\delta]\) is their common successor.
Thus the quotient is globally convergent, hence locally directed.
The quotient size is still at most \(2^s\); generating did not
change the sentence set. This supplies the full S4.2 construction,
not a claim that arbitrary locally directed filtrations preserve G.

In every case we have a finite frame in the correct class with a
world falsifying A. Conversely, every theorem is true on every
finite frame in the class by soundness. Thus finite models determine
theoremhood and give the asserted bound. This reasoning constructs
one finite model for one finite sentence; it does not replace the
infinite canonical model for an arbitrary entire premise set.
\end{proof}
\end{annexitem}

\begin{annexitem}{annex:p8.chellas.decision}{thm:p08-chellas-decision}{Finite search termination, sound answers and source dovetailing}
\begin{proof}
For j worlds there are \(2^{j^2}\) binary relations and, if A has
t distinct atoms, \(2^{jt}\) relevant valuations. Each relation
test is a finite Boolean computation on its adjacency matrix:
seriality checks a successor in each row, reflexivity checks the
diagonal, symmetry checks reversed entries, transitivity and
Euclideanness check each triple, and directedness checks a common
successor for each pair of successors of each world. Every retained
frame satisfies the exact chosen class, by these definitions.
Atomic valuation for atoms not in A can be fixed arbitrarily,
since structural induction shows that they cannot affect its truth.
Bottom-up truth evaluation assigns each subsentence a bit at each
world. Boolean cases are truth-table operations; each modal case
scans a finite successor row, with empty rows interpreted
universally or existentially as stipulated. Thus each test terminates.

The bound b is an effectively computed positive integer from the
finite parse tree. There are finitely many j-values, relations,
valuations and worlds to test, so the entire search terminates.
A discovered countermodel proves non-theoremhood by soundness.
If no countermodel is found, the bound in the previous theorem
excludes non-theoremhood: every countermodel of size at most b
is isomorphic to a labeled candidate in the enumeration, and
truth is preserved by the isomorphism by structural induction.
Hence the returned theorem answer is correct too.

For a finite premise list, local consequence is equivalent by
the truth clauses to validity of its conjunction-implies-A
sentence. The source consequence definition and deduction theorem
give the same syntactic reduction, including the empty list.
Satisfiability of A is equivalent to failure of validity of
\(\neg A\); completeness converts that to non-theoremhood of
\(\neg A\). These prove the two further decisions.

For the source's alternative procedure, enumerate finite proof
sequences with their rule annotations and separate no-premise
proofs for necessitation. Recognition is effective: finite formula
parsing, tautological consequence by truth tables on propositional
atoms, schema matching, K/dual schema matching and rule checking
are all finite tests. Therefore every theorem eventually appears
as a proved final sentence. Independently enumerate all finite
labeled frames of every size and relevant valuations, with the
same effective frame and truth tests. Every non-theorem eventually
appears as a countermodel by finite determination. Alternate one
bounded step from each search. Soundness prevents opposite valid
answers, and one search must succeed, proving termination even
without using the explicit size bound. Neither proof enumeration
alone nor unbounded model enumeration alone decides both outcomes.
There is no finite-input representation or terminating search
asserted here for an arbitrary infinite premise set.
\end{proof}
\end{annexitem}

\begin{annexitem}{annex:p8.chellas.modalities}{thm:p08-chellas-modalities}{Reduction laws, exact finite witnesses and infinite families}
\begin{proof}
Push negations in a modality to its outside by
\(\Box\neg A\leftrightarrow\neg\Diamond A\) and
\(\Diamond\neg A\leftrightarrow\neg\Box A\), and eliminate
double negations. Every modality is consequently a box/diamond
word, possibly preceded by one negation.

In S4, T and 4 give \(\Box\Box A\leftrightarrow\Box A\),
and duality gives \(\Diamond\Diamond A\leftrightarrow\Diamond A\).
It also has
\(\Diamond\Box A\leftrightarrow
\Diamond\Box\Diamond\Box A\).
For the forward implication, the T-dual instance yields
\(\Box A\to\Diamond\Box A\); box monotonicity and 4
give \(\Box A\to\Box\Diamond\Box A\); diamond
monotonicity gives the desired implication. For the reverse,
T gives \(\Box\Diamond\Box A\to\Diamond\Box A\);
diamond monotonicity gives a double-diamond consequent,
which collapses by dual 4. The dual identity is
\(\Box\Diamond A\leftrightarrow
\Box\Diamond\Box\Diamond A\).
First remove adjacent equal operators. Any remaining word
alternates, and a length-four alternating block reduces to its
first two operators by these two identities. Repetition reduces
length to at most three, giving precisely the seven affirmative
representatives in the statement as an upper bound.

In S4.2, 4 gives
\(\Diamond\Box A\to\Diamond\Box\Box A\).
G applied to \(\Box A\) then gives
\(\Diamond\Box A\to\Box\Diamond\Box A\).
The reverse follows from T. Hence
\(\Box\Diamond\Box A\leftrightarrow\Diamond\Box A\),
and duality gives
\(\Diamond\Box\Diamond A\leftrightarrow\Box\Diamond A\).
These reduce the S4 seven to the stated five. In S5, 5 and T
give \(\Box\Diamond A\leftrightarrow\Diamond A\), and
the dual gives \(\Diamond\Box A\leftrightarrow\Box A\).
Together with the repeated-operator reductions this leaves the
stated three affirmative modalities.

We show these upper bounds exact with a finite truth matrix.
Let F have worlds a,b,c, all reflexive edges and only extra
edges \((a,b),(a,c)\); let H have the same worlds, all reflexive
edges, extra edges \((a,b),(a,c)\), and both directions between
b and c. F is S4, and H is globally convergent S4.2. The following
are truth values at a, with precisely the indicated valuation of P.
\begin{center}\small
\begin{tabular}{@{}llrrrrrrr@{}}
\toprule
Frame & P & identity & \(\Box\) & \(\Diamond\) &
 \(\Box\Diamond\) & \(\Diamond\Box\) &
 \(\Box\Diamond\Box\) & \(\Diamond\Box\Diamond\) \\
\midrule
F & \(\{a\}\) & 1&0&1&0&0&0&0 \\
F & \(\{b\}\) & 0&0&1&0&1&0&1 \\
F & \(\{b,c\}\) & 0&0&1&1&1&1&1 \\
H & \(\{b\}\) & 0&0&1&1&0&0&1 \\
\bottomrule
\end{tabular}
\end{center}
These values follow by scanning successors: in F the terminal
worlds see themselves only, while in H each terminal world sees
both b and c. All seven columns are different, so every pair
has a finite S4 counterexample to its alleged equivalence.
For the S4.2 five, use H with valuations
\(\{a\},\{b\},\{b,c\}\). Their first five columns are,
respectively,
\((1,0,1,0,0),(0,0,1,1,0),(0,0,1,1,1)\),
again distinct column by column. For S5, use the two-world
universal frame with P true at a alone. Evaluating at a and b
gives identity/box/diamond rows \((1,0,1),(0,0,1)\),
so the three columns are distinct. For each finite count,
negating distinct columns preserves distinctness. No affirmative
representative equals any negative one: on the one-world reflexive
frame with P false, all affirmative representatives are false
and every negative one true. The exact counts are therefore
fourteen, ten and six.

For the infinite counts in K, KD, KT, KB and KTB, fix \(n>m\geq0\)
and take the reflexive symmetric path with worlds
\(a_0,\ldots,a_n\), where exactly equal or neighboring indices
are related. Put P true at \(a_0,\ldots,a_m\).
At \(a_0\), every m-step reachable world has index at most m,
so \(\Box^m P\) is true. The n-step forward path reaches
\(a_n\), where P is false, so \(\Box^n P\) is false.
This model belongs to KTB and to each of its selected subsystems,
including KB. It separates every two iterated boxes in those five
systems, proving infinitude without a reduction bound assumption.

For K4 take worlds \(0,\ldots,n\), put \(iRj\) when \(i<j\)
or \(i=j=n\), and make P true only at n. The relation is serial
and transitive. Every n-step path from 0 ends at n, because
each step before n strictly increases its index, while n can
loop. Thus \(\Box^n P\) is true. The m-step path through
\(0,1,\ldots,m\) ends below n and falsifies P, so
\(\Box^m P\) is false; when m=0 it is simply false at 0.
This separates all iterated boxes in K4. Together with the
preceding five cases, all six infinite counts follow.
\end{proof}
\end{annexitem}
